% ===== VARIANT: PhD APPLICATION (QUANTITATIVE FINANCE) =============
% Academic CV. Research interests and research experience lead;
% teaching, awards and references retained. Two pages acceptable.
% ===================================================================
% =====================================================================
%  Shared preamble for all CV variants.
%  Each variant does:   % =====================================================================
%  Shared preamble for all CV variants.
%  Each variant does:   % =====================================================================
%  Shared preamble for all CV variants.
%  Each variant does:   \input{cv-preamble.tex}
%  Do not put content here — only packages, formatting and macros.
% =====================================================================
\documentclass[letterpaper,11pt]{article}

\usepackage{latexsym}
\usepackage[empty]{fullpage}
\usepackage{titlesec}
\usepackage{marvosym}
\usepackage[usenames,dvipsnames]{color}
\usepackage{verbatim}
\usepackage{enumitem}
\usepackage[hidelinks]{hyperref}
\usepackage{fancyhdr}
\usepackage[english]{babel}
\usepackage{tabularx}
\usepackage{fontawesome5}
\usepackage{multicol}
\usepackage{graphicx}
\setlength{\multicolsep}{-3.0pt}
\setlength{\columnsep}{-1pt}
\input{glyphtounicode}
\usepackage{paracol}

\pagestyle{fancy}
\fancyhf{}
\fancyfoot{}
\renewcommand{\headrulewidth}{0pt}
\renewcommand{\footrulewidth}{0pt}

\addtolength{\oddsidemargin}{-0.5in}
\addtolength{\evensidemargin}{-0.5in}
\addtolength{\textwidth}{0.9in}
\addtolength{\topmargin}{-.3in}
\addtolength{\textheight}{0.9in}
\hypersetup{
    colorlinks=true,
    linkcolor=blue,
    urlcolor=blue,
    filecolor=blue,
    pdfpagemode=FullScreen
}
\urlstyle{same}

\raggedbottom
\raggedright
\setlength{\tabcolsep}{0in}

\titleformat{\section}{
  \vspace{-4pt}\scshape\raggedright\large\bfseries
}{}{0em}{}[\color{black}\titlerule \vspace{-5pt}]

\pdfgentounicode=1

\newcommand{\resumeItem}[1]{
  \item\small{{#1 \vspace{-2pt}}}
}

\newcommand{\resumeSubheading}[4]{
  \vspace{-2pt}\item
    \begin{tabular*}{1.0\textwidth}[t]{l@{\extracolsep{\fill}}r}
      \textbf{#1} & \textbf{\small #2} \\
      \textit{\small#3} & \textit{\small #4} \\
    \end{tabular*}\vspace{-7pt}
}

\newcommand{\resumeProjectHeading}[2]{
    \item
    \begin{tabular*}{0.97\textwidth}{l@{\extracolsep{\fill}}r}
      \small#1 & \textbf{\small #2}\\
    \end{tabular*}\vspace{-7pt}
}

\renewcommand\labelitemi{$\vcenter{\hbox{\tiny$\bullet$}}$}
\renewcommand\labelitemii{$\vcenter{\hbox{\tiny$\bullet$}}$}

\newcommand{\resumeSubHeadingListStart}{\begin{itemize}[leftmargin=0.0in, label={}]}
\newcommand{\resumeSubHeadingListEnd}{\end{itemize}}
\newcommand{\resumeItemListStart}{\begin{itemize}}
\newcommand{\resumeItemListEnd}{\end{itemize}\vspace{-5pt}}

% ---------------------------------------------------------------------
% \cvheader{positioning line}{focus line}
% ---------------------------------------------------------------------
\newcommand{\cvheader}[2]{%
\begin{center}
    {\Huge \scshape Franck Deba} \\[0.4em]
    {\small \textit{#1}} \\[0.25em]
    {\small #2}
\end{center}
\vspace{0.5pt}
\small
\raisebox{-0.1\height}{\faPhone}~+33 (0)6 05 77 47 51 \hspace{1em}
\raisebox{-0.2\height}{\faEnvelope}~\href{mailto:debafranck@gmail.com}{\underline{debafranck@gmail.com}} \hspace{1em}
\raisebox{-0.2\height}{\faLinkedin}~\href{https://linkedin.com/in/deba-franck}{\underline{deba-franck}} \hspace{1em}
\raisebox{-0.2\height}{\faGithub}~\href{https://github.com/usenameFD}{\underline{usenameFD}} \hspace{1em}
\raisebox{-0.2\height}{\faGlobe}~\href{https://usenamefd.github.io/}{\underline{usenamefd.github.io}}
\vspace{-8pt}
}

%  Do not put content here — only packages, formatting and macros.
% =====================================================================
\documentclass[letterpaper,11pt]{article}

\usepackage{latexsym}
\usepackage[empty]{fullpage}
\usepackage{titlesec}
\usepackage{marvosym}
\usepackage[usenames,dvipsnames]{color}
\usepackage{verbatim}
\usepackage{enumitem}
\usepackage[hidelinks]{hyperref}
\usepackage{fancyhdr}
\usepackage[english]{babel}
\usepackage{tabularx}
\usepackage{fontawesome5}
\usepackage{multicol}
\usepackage{graphicx}
\setlength{\multicolsep}{-3.0pt}
\setlength{\columnsep}{-1pt}
\input{glyphtounicode}
\usepackage{paracol}

\pagestyle{fancy}
\fancyhf{}
\fancyfoot{}
\renewcommand{\headrulewidth}{0pt}
\renewcommand{\footrulewidth}{0pt}

\addtolength{\oddsidemargin}{-0.5in}
\addtolength{\evensidemargin}{-0.5in}
\addtolength{\textwidth}{0.9in}
\addtolength{\topmargin}{-.3in}
\addtolength{\textheight}{0.9in}
\hypersetup{
    colorlinks=true,
    linkcolor=blue,
    urlcolor=blue,
    filecolor=blue,
    pdfpagemode=FullScreen
}
\urlstyle{same}

\raggedbottom
\raggedright
\setlength{\tabcolsep}{0in}

\titleformat{\section}{
  \vspace{-4pt}\scshape\raggedright\large\bfseries
}{}{0em}{}[\color{black}\titlerule \vspace{-5pt}]

\pdfgentounicode=1

\newcommand{\resumeItem}[1]{
  \item\small{{#1 \vspace{-2pt}}}
}

\newcommand{\resumeSubheading}[4]{
  \vspace{-2pt}\item
    \begin{tabular*}{1.0\textwidth}[t]{l@{\extracolsep{\fill}}r}
      \textbf{#1} & \textbf{\small #2} \\
      \textit{\small#3} & \textit{\small #4} \\
    \end{tabular*}\vspace{-7pt}
}

\newcommand{\resumeProjectHeading}[2]{
    \item
    \begin{tabular*}{0.97\textwidth}{l@{\extracolsep{\fill}}r}
      \small#1 & \textbf{\small #2}\\
    \end{tabular*}\vspace{-7pt}
}

\renewcommand\labelitemi{$\vcenter{\hbox{\tiny$\bullet$}}$}
\renewcommand\labelitemii{$\vcenter{\hbox{\tiny$\bullet$}}$}

\newcommand{\resumeSubHeadingListStart}{\begin{itemize}[leftmargin=0.0in, label={}]}
\newcommand{\resumeSubHeadingListEnd}{\end{itemize}}
\newcommand{\resumeItemListStart}{\begin{itemize}}
\newcommand{\resumeItemListEnd}{\end{itemize}\vspace{-5pt}}

% ---------------------------------------------------------------------
% \cvheader{positioning line}{focus line}
% ---------------------------------------------------------------------
\newcommand{\cvheader}[2]{%
\begin{center}
    {\Huge \scshape Franck Deba} \\[0.4em]
    {\small \textit{#1}} \\[0.25em]
    {\small #2}
\end{center}
\vspace{0.5pt}
\small
\raisebox{-0.1\height}{\faPhone}~+33 (0)6 05 77 47 51 \hspace{1em}
\raisebox{-0.2\height}{\faEnvelope}~\href{mailto:debafranck@gmail.com}{\underline{debafranck@gmail.com}} \hspace{1em}
\raisebox{-0.2\height}{\faLinkedin}~\href{https://linkedin.com/in/deba-franck}{\underline{deba-franck}} \hspace{1em}
\raisebox{-0.2\height}{\faGithub}~\href{https://github.com/usenameFD}{\underline{usenameFD}} \hspace{1em}
\raisebox{-0.2\height}{\faGlobe}~\href{https://usenamefd.github.io/}{\underline{usenamefd.github.io}}
\vspace{-8pt}
}

%  Do not put content here — only packages, formatting and macros.
% =====================================================================
\documentclass[letterpaper,11pt]{article}

\usepackage{latexsym}
\usepackage[empty]{fullpage}
\usepackage{titlesec}
\usepackage{marvosym}
\usepackage[usenames,dvipsnames]{color}
\usepackage{verbatim}
\usepackage{enumitem}
\usepackage[hidelinks]{hyperref}
\usepackage{fancyhdr}
\usepackage[english]{babel}
\usepackage{tabularx}
\usepackage{fontawesome5}
\usepackage{multicol}
\usepackage{graphicx}
\setlength{\multicolsep}{-3.0pt}
\setlength{\columnsep}{-1pt}
\input{glyphtounicode}
\usepackage{paracol}

\pagestyle{fancy}
\fancyhf{}
\fancyfoot{}
\renewcommand{\headrulewidth}{0pt}
\renewcommand{\footrulewidth}{0pt}

\addtolength{\oddsidemargin}{-0.5in}
\addtolength{\evensidemargin}{-0.5in}
\addtolength{\textwidth}{0.9in}
\addtolength{\topmargin}{-.3in}
\addtolength{\textheight}{0.9in}
\hypersetup{
    colorlinks=true,
    linkcolor=blue,
    urlcolor=blue,
    filecolor=blue,
    pdfpagemode=FullScreen
}
\urlstyle{same}

\raggedbottom
\raggedright
\setlength{\tabcolsep}{0in}

\titleformat{\section}{
  \vspace{-4pt}\scshape\raggedright\large\bfseries
}{}{0em}{}[\color{black}\titlerule \vspace{-5pt}]

\pdfgentounicode=1

\newcommand{\resumeItem}[1]{
  \item\small{{#1 \vspace{-2pt}}}
}

\newcommand{\resumeSubheading}[4]{
  \vspace{-2pt}\item
    \begin{tabular*}{1.0\textwidth}[t]{l@{\extracolsep{\fill}}r}
      \textbf{#1} & \textbf{\small #2} \\
      \textit{\small#3} & \textit{\small #4} \\
    \end{tabular*}\vspace{-7pt}
}

\newcommand{\resumeProjectHeading}[2]{
    \item
    \begin{tabular*}{0.97\textwidth}{l@{\extracolsep{\fill}}r}
      \small#1 & \textbf{\small #2}\\
    \end{tabular*}\vspace{-7pt}
}

\renewcommand\labelitemi{$\vcenter{\hbox{\tiny$\bullet$}}$}
\renewcommand\labelitemii{$\vcenter{\hbox{\tiny$\bullet$}}$}

\newcommand{\resumeSubHeadingListStart}{\begin{itemize}[leftmargin=0.0in, label={}]}
\newcommand{\resumeSubHeadingListEnd}{\end{itemize}}
\newcommand{\resumeItemListStart}{\begin{itemize}}
\newcommand{\resumeItemListEnd}{\end{itemize}\vspace{-5pt}}

% ---------------------------------------------------------------------
% \cvheader{positioning line}{focus line}
% ---------------------------------------------------------------------
\newcommand{\cvheader}[2]{%
\begin{center}
    {\Huge \scshape Franck Deba} \\[0.4em]
    {\small \textit{#1}} \\[0.25em]
    {\small #2}
\end{center}
\vspace{0.5pt}
\small
\raisebox{-0.1\height}{\faPhone}~+33 (0)6 05 77 47 51 \hspace{1em}
\raisebox{-0.2\height}{\faEnvelope}~\href{mailto:debafranck@gmail.com}{\underline{debafranck@gmail.com}} \hspace{1em}
\raisebox{-0.2\height}{\faLinkedin}~\href{https://linkedin.com/in/deba-franck}{\underline{deba-franck}} \hspace{1em}
\raisebox{-0.2\height}{\faGithub}~\href{https://github.com/usenameFD}{\underline{usenameFD}} \hspace{1em}
\raisebox{-0.2\height}{\faGlobe}~\href{https://usenamefd.github.io/}{\underline{usenamefd.github.io}}
\vspace{-8pt}
}


\begin{document}

\cvheader{PhD Candidate --- Quantitative Finance}{Stochastic control \& optimal stopping $\cdot$ Filtering and latent-factor calibration $\cdot$ Numerical methods for PDEs $\cdot$ Risk measures}

%-----------RESEARCH INTERESTS-----------
\section{Research Interests}
\begin{itemize}[leftmargin=0.15in, label={}]
  \small{\item{
    Stochastic optimal control and Hamilton--Jacobi--Bellman(--Isaacs) equations; optimal stopping and free-boundary problems;
    calibration and identifiability of latent factor models; sequential Monte Carlo and filtering for stochastic volatility;
    dependence structures and extreme value theory for portfolio credit risk; numerical analysis of PDEs and
    neural-network-based solvers.
  }}
\end{itemize}

%-----------EDUCATION-----------
\section{Education}
\resumeSubHeadingListStart
  \resumeSubheading
    {Master M2MO --- Mod\'elisation Al\'eatoire (Quantitative Finance and Data Science)}{Sep 2025 -- Present}
    {Universit\'e Paris Cit\'e}{Paris, France}
  \vspace{3pt}
  \resumeSubheading
    {MSc in Data Science and Risk Management}{Aug 2023 -- Dec 2025}
    {ENSAI (\'Ecole Nationale de la Statistique et de l'Analyse de l'Information)}{Rennes, France}
  \vspace{3pt}
  \resumeSubheading
    {MSc in Statistics applied to Economics}{Sep 2021 -- Aug 2023}
    {ISSEA -- CEMAC}{Yaound\'e, Cameroon}
  \vspace{3pt}
  \resumeSubheading
    {Bachelor's Degree in Mathematics}{Sep 2017 -- Aug 2020}
    {Double degree --- Higher Teacher Training College \& University of Yaound\'e I}{Yaound\'e, Cameroon}
\resumeSubHeadingListEnd

%-----------MASTER'S THESIS-----------
\section{Master's Thesis}
\resumeSubHeadingListStart
  \resumeProjectHeading
    {\href{https://usenamefd.github.io/projects/drc-frtb-cac40.html}{\textbf{Factor-Model Calibration for the Default Risk Charge}} $|$ \emph{47 pp., Universit\'e Paris Cit\'e}} {2026}
    \resumeItemListStart
      \resumeItem{\textbf{Question.} Part of the DRC literature implicitly assumes the calibration method is secondary once a Gaussian multi-factor framework is correctly posed. The thesis tests that assumption.}
      \resumeItem{\textbf{Contribution.} Introduces calibration by maximum likelihood of the model --- to my knowledge a new contribution to DRC modelling --- alongside iterative PCA and Frobenius-norm minimisation, on a portfolio of 37 issuers.}
      \resumeItem{\textbf{Theory.} Establishes that the likelihood criterion is equivalent to minimising a Kullback--Leibler divergence, whose penalisation geometry is asymmetric: it penalises correlation under-estimation more severely where the target correlation is already high. Discusses non-compactness of the likelihood programme, Heywood cases, and non-identifiability of loadings up to rotation.}
      \resumeItem{\textbf{Empirics.} Iterative PCA and Frobenius converge to an identical fit; maximum likelihood departs by under 3\% in relative Frobenius error, concentrated almost entirely on the financial sub-group. Marchenko--Pastur independently selects two factors on all three calibration windows.}
      \resumeItem{\textbf{Inference.} Monte Carlo with 200{,}000 scenarios and confidence intervals by sectioning (Dong \& Nakayama, 2020), appropriate for a step-valued loss distribution.}
      \resumeItem{\textbf{Diagnostics.} Systematic stationarity (ADF, KPSS) and normality (Shapiro--Wilk, Kolmogorov--Smirnov) testing of factors and residuals, addressing a gap in the reviewed literature.}
    \resumeItemListEnd
\resumeSubHeadingListEnd

%-----------RESEARCH PROJECTS-----------
\section{Research Projects}
\resumeSubHeadingListStart

  \resumeProjectHeading
    {\href{https://github.com/usenameFD/Conjecture-de-Collatz}{\textbf{An Analytic Reformulation of the Collatz Conjecture}} $|$ \emph{Markov chains}} {2026}
    \resumeItemListStart
      \resumeItem{Reformulated the conjecture as a uniqueness problem for the stationary measure of a Markov chain}
      \resumeItem{Characterised admissible stationary measures and verified the characterisation numerically}
    \resumeItemListEnd

  \resumeProjectHeading
    {\href{https://github.com/usenameFD/Numerical_methods/blob/main/zermeloNotebook.ipynb}{\textbf{Numerical Resolution of an HJBI PDE --- Zermelo Stochastic Control Problem}} $|$ \emph{Python, PyTorch}} {2026}
    \resumeItemListStart
      \resumeItem{Solved a Hamilton--Jacobi--Bellman--Isaacs PDE by finite differences with a semi-smooth Newton method, the equation being non-linear and non-smooth}
      \resumeItem{Measured the numerical order of convergence against a known closed-form solution, then applied the scheme to a case with no explicit solution}
      \resumeItem{Developed a Physics-Informed Neural Network solver using automatic differentiation (Esteve-Yag\"ue, 2025); compared three loss formulations and studied sensitivity to the loss weights}
      \resumeItem{Concluded on the trade-off: provable convergence order for the deterministic scheme against dimensional scalability for the network}
    \resumeItemListEnd

  \resumeProjectHeading
    {\href{https://usenamefd.github.io/projects/methodes-numeriques-edp.html}{\textbf{Optimal Stopping and Convergence to the Brownian Bridge}} $|$ \emph{Python}} {2026}
    \resumeItemListStart
      \resumeItem{Modelled a card-game brainteaser (Crack, 2014) as an optimal stopping problem solved by backward induction}
      \resumeItem{Computed the discrete stopping boundary and its convergence to the continuous boundary; asymptotic analysis of the normalised value $V(N,N)/\sqrt{2N}$ following Ekstr\"om \& Vaicenavicius (2020)}
    \resumeItemListEnd

  \resumeProjectHeading
    {\href{https://usenamefd.github.io/projects/calibration-heston-filtrage.html}{\textbf{Sequential Monte Carlo for Stochastic Volatility Calibration}} $|$ \emph{Python, R}} {2025}
    \resumeItemListStart
      \resumeItem{Heston model written as a state-space system with the exact non-central chi-square transition law of the CIR variance}
      \resumeItem{Bootstrap and auxiliary particle filters implemented and compared in a simulation study across option moneyness, where the true latent trajectory is known}
      \resumeItem{Companion study on the Taylor log-SV model: Kalman filtering after linearisation vs.\ bootstrap particle filtering, quantifying the error of the Gaussian approximation of a log-$\chi^2_1$ observation noise}
    \resumeItemListEnd

  \resumeProjectHeading
    {\href{https://usenamefd.github.io/projects/creditvar-copules.html}{\textbf{Dependence Structures for Portfolio Credit Risk}} $|$ \emph{Python}} {2025}
    \resumeItemListStart
      \resumeItem{Two-step estimation (Sklar): skewed-Student margins by maximum likelihood, then elliptical and Archimedean copulas}
      \resumeItem{Empirical tail-dependence evidence incompatible with the asymptotic independence implied by the Gaussian copula, with consequences for capital}
    \resumeItemListEnd

  \resumeProjectHeading
    {\href{https://usenamefd.github.io/projects/value-at-risk.html}{\textbf{Extreme Value Theory for Risk Measures}} $|$ \emph{Python, R}} {2025}
    \resumeItemListStart
      \resumeItem{Block maxima (GEV vs.\ Gumbel, likelihood-ratio testing) and peaks-over-threshold (GPD) with threshold calibration and goodness-of-fit testing}
      \resumeItem{GARCH-filtered residual approach to relax the i.i.d.\ assumption; backtesting by unconditional-coverage testing}
    \resumeItemListEnd

\resumeSubHeadingListEnd

%-----------RESEARCH & PROFESSIONAL EXPERIENCE-----------
\section{Research \& Professional Experience}
\resumeSubHeadingListStart
  \resumeSubheading
    {\textbf{Quantitative Analyst (Intern)} $|$ \textbf{PwC}}{Apr -- Oct 2026}
    {Risk Line of Service}{Paris, France}
    \resumeItemListStart
      \resumeItem{Research internship supporting the Master's thesis on Default Risk Charge modelling under FRTB}
    \resumeItemListEnd

  \resumeSubheading
    {\textbf{Quantitative Risk Analyst (Intern)} $|$ \textbf{Cr\'edit Agricole}}{Apr -- Sep 2025}
    {Risk Modeling Department}{Paris, France}
    \resumeItemListStart
      \resumeItem{Climate stress-testing methodologies using time series and Monte Carlo methods}
      \resumeItem{Automated backtesting tool in Python (Dash) for model validation}
    \resumeItemListEnd

  \resumeSubheading
    {\textbf{Data Science Intern} $|$ \textbf{Groupe VYV -- Harmonie Mutuelle}}{Jun -- Aug 2024}
    {Assurance Pr\'evoyance}{Paris, France}
    \resumeItemListStart
      \resumeItem{Survival analysis (Cox model) and clustering applied to insurance risk profiling}
    \resumeItemListEnd
\resumeSubHeadingListEnd

%-----------TEACHING-----------
\section{Teaching Experience}
\resumeSubHeadingListStart
  \resumeSubheading
    {\textbf{Intern Mathematics Teacher} $|$ \textbf{Ngoa Ekel\'e High School}}{Jan -- May 2022}
    {Mathematics Department}{Yaound\'e, Cameroon}
    \resumeItemListStart
      \resumeItem{Taught mathematics using a competency-based approach}
      \resumeItem{Designed assessments, analysed student performance and developed targeted remediation strategies}
    \resumeItemListEnd
\resumeSubHeadingListEnd

%-----------DISSEMINATION-----------
\section{Dissemination}
\begin{itemize}[leftmargin=0.15in, label={}]
  \small{\item{
    Maintains \href{https://usenamefd.github.io/}{\underline{usenamefd.github.io}}, a documented research portfolio presenting the above work with its
    mathematical framework, data provenance, assumptions and stated limitations, together with open-source implementations.
  }}
\end{itemize}

%-----------SKILLS-----------
\section{Technical Skills}
\begin{itemize}[leftmargin=0.15in, label={}]
  \small{\item{
    \textbf{Languages:} Python (NumPy, SciPy, PyTorch, scikit-learn, statsmodels), R, C++ (basics), SQL, SAS, \LaTeX \\
    \textbf{Methods:} Stochastic control, HJB/HJBI PDEs, optimal stopping, finite differences (Euler, Crank--Nicolson, PSOR, BDF), semi-smooth Newton, PINNs, sequential Monte Carlo, Kalman filtering, EM, spectral projected gradient, random matrix theory, copulas, extreme value theory
  }}
\end{itemize}

%-----------AWARDS & LANGUAGES-----------
\section{Awards, Certifications \& Languages}
\begin{itemize}[leftmargin=0.15in, label={}]
  \small{\item{
    \textbf{Awards:} National Essay Competition Prize (2019 --- Ministry of External Relations, Cameroon), PK Fokam Prize for Science Students (2017) \\
    \textbf{Certifications:} AMF (Autorit\'e des March\'es Financiers), Financial Markets (Yale), AI in Finance (World Bank) \\
    \textbf{Languages:} French (native), English (advanced, B2 --- TOEIC/TOEFL)
  }}
\end{itemize}

%-----------REFERENCES-----------
\section{References}
\begin{itemize}[leftmargin=0.15in, label={}]
  \small{\item{
    \textbf{Adrien Saumard} --- ENSAI, Rennes \\
    \textbf{Basile De Loynes} --- ENSAI, Rennes \\
    \textbf{Sim\'eon Fotso} --- \'Ecole Normale Sup\'erieure, Yaound\'e
  }}
\end{itemize}

\end{document}
